% !TeX root = ../../Book5.tex
% 纲要标题：### 21.1 余代数与 Hopf 代数
% 纲要位置：工作记录/计划纲要.md，第 4391 行。

\section{余代数与 Hopf 代数}
\label{b5:sec:2101}


% 纲要范围：
%
% * 余乘法、余单位与对极
% * 群代数及坐标环的对偶方向
% * 模与余模的区别
%

群表示与 Lie 代数表示都能作张量积，但作用公式不同：群元素同时作用于两个因子，Lie 代数元素则分别作用后相加。余乘法把这两种规则统一写成代数中的映射；再把单位表示和对偶表示所需的规则一并加入，就得到 Hopf 代数。

\subsection{余乘法、余单位与对极}
\label{b5:subsec:2101:01}

\begin{definition}\label{b5:def:coalgebra-hopf}
设 \(k\) 为域，\(H\) 为 \(k\) 向量空间。线性映射
\(\Delta:H\to H\otimes_kH\)、\(\varepsilon:H\to k\) 满足
\[
(\Delta\otimes1)\Delta=(1\otimes\Delta)\Delta,\qquad
(\varepsilon\otimes1)\Delta=1_H=(1\otimes\varepsilon)\Delta
\]
时，称 \((H,\Delta,\varepsilon)\) 为\term[yudaishu]{余代数}[coalgebra]，称两映射为余乘法与余单位；后两式用 \(k\otimes H=H=H\otimes k\) 识别。
若 \(H\) 还为含幺结合 \(k\) 代数，且 \(\Delta,\varepsilon\) 是含幺代数同态，就称为\term[shuangdaishu]{双代数}[bialgebra]。若另有线性映射 \(S:H\to H\) 满足
\[
m(S\otimes1)\Delta=\eta\varepsilon=m(1\otimes S)\Delta,
\]
其中 \(m\) 是乘法、\(\eta:k\to H\) 是单位映射，就称为\term[Hopf-daishu]{Hopf 代数}[Hopf algebra]，称 \(S\) 为\term[duiji]{对极}[antipode]。
\end{definition}
\symidx{\(\Delta,\varepsilon,S\)：Hopf 代数的余乘法、余单位与对极}
本书不把对极的可逆性额外放进定义；它在本章的群代数、包络代数及交换坐标代数例子中成立。Etingof 等的《Tensor Categories》在定义 Hopf 代数时要求对极可逆，因此引用其张量范畴结论时须保留这一差别。
\ProseParagraphBreak
常把 \(\Delta(h)\) 写为 \(\sum h_{(1)}\otimes h_{(2)}\)，称为 Sweedler 记号。它表示一个有限和，不表示每个 \(h\) 都只有一个纯张量分解。余结合律允许把三重余乘法统一写为 \(\sum h_{(1)}\otimes h_{(2)}\otimes h_{(3)}\)。
\symidx{\(\sum h_{(1)}\otimes h_{(2)}\)：余乘法的 Sweedler 记号}
\begin{proposition}\label{b5:prop:antipode-unique}
设 \(k\) 为域，\(H\) 为 \(k\) 双代数。其对极 \(S:H\to H\) 若存在，则唯一，并满足 \(S(1)=1\)、\(S(ab)=S(b)S(a)\) 对全部 \(a,b\in H\) 成立。
\end{proposition}
\begin{proof}
在 \(\operatorname{Hom}_k(H,H)\) 上定义卷积
\((f*g)(h)=\sum f\bigl(h_{(1)}\bigr)g\bigl(h_{(2)}\bigr)\)。
结合律来自 \(m\) 的结合与 \(\Delta\) 的余结合，单位为 \(\eta\varepsilon\)。对极就是恒等映射的双边卷积逆，故唯一；代入单位即得 \(S(1)=1\)。
对反乘法性质，改在 \(\operatorname{Hom}_k(H\otimes H,H)\) 中卷积，张量积余乘法把 \(a\otimes b\) 送到
\(\sum\bigl(a_{(1)}\otimes b_{(1)}\bigr)\otimes\bigl(a_{(2)}\otimes b_{(2)}\bigr)\)。
映射 \(S\circ m\) 是乘法映射 \(m\) 的卷积逆，因 \(\Delta\) 保持乘法。另一个映射 \(F(a,b)=S(b)S(a)\) 也给出逆：两侧复合分别为
\[
\sum S\bigl(b_{(1)}\bigr)S\bigl(a_{(1)}\bigr)a_{(2)}b_{(2)}
=\varepsilon(a)\varepsilon(b)1
\]
和
\(\sum a_{(1)}b_{(1)}S\bigl(b_{(2)}\bigr)S\bigl(a_{(2)}\bigr)=\varepsilon(a)\varepsilon(b)1\)。
逆的唯一性使两映射相同。
\end{proof}
若再有交换 \(\Delta\) 两因子不改变它，则称余交换。乘法交换与余乘法余交换是两项不同条件，以下两种来自同一个群的代数将说明这种区别。

\subsection{群代数及坐标环的对偶方向}
\label{b5:subsec:2101:02}

\begin{proposition}\label{b5:prop:group-enveloping-hopf}
设 \(k\) 为域，\(G\) 为群，\(\mathfrak g\) 为 \(k\) 上 Lie 代数。下列公式分别给出 \(kG\) 与 \(U(\mathfrak g)\) 上的 Hopf 代数结构：
\[
\begin{array}{c|ccc}
&H:\Delta&\varepsilon&S\\ \hline
kG&g\mapsto g\otimes g&g\mapsto1&g\mapsto g^{-1}\\
U(\mathfrak g)&x\mapsto x\otimes1+1\otimes x&x\mapsto0&x\mapsto-x
\end{array}
\]
第二行的公式在 \(x\in\mathfrak g\) 上规定，再按相应代数或反代数映射延拓。两者都余交换。
\end{proposition}
\begin{proof}
群代数的各式在群基元上核验公理即可，例如
\(m(S\otimes1)\Delta(g)=g^{-1}g=1=\varepsilon(g)1\)。
包络代数中，映射 \(x\mapsto x\otimes1+1\otimes x\) 保持 Lie 括号，因为交叉的两个因子交换，故由泛性质延拓为代数同态。增广零化全部 Lie 括号；映射 \(x\mapsto-x\) 到反代数保持 Lie 括号，因而给出反同态 \(S\)。余结合及余单位等式在生成元上成立，所以在全部代数上成立。对极等式可对字长归纳：若两较短字 \(u,v\) 满足左式，则
\[
\sum S\bigl((uv)_{(1)}\bigr)(uv)_{(2)}
=\sum S\bigl(v_{(1)}\bigr)S\bigl(u_{(1)}\bigr)u_{(2)}v_{(2)}
=\varepsilon(u)\varepsilon(v)1.
\]
右式同理；线性性完成证明。
\end{proof}
若 \(G\) 有限，函数空间 \(k^G\) 则按逐点乘法成为交换 Hopf 代数：
\[
\Delta f(g,h)=f(gh),\qquad \varepsilon(f)=f(1),\qquad
Sf(g)=f(g^{-1}).
\]
在点函数基 \(\delta_x\) 中，
\(\Delta\delta_x=\sum_{ab=x}\delta_a\otimes\delta_b\)。
这是 \(kG\) 的线性对偶 Hopf 代数；函数代数通常不余交换，恰在 \(G\) 交换时余交换。群无限时，任意函数的两变量拉回未必在代数张量积 \(k^G\otimes k^G\) 中，因此不能不加限制照搬有限对偶。

\ProseParagraphBreak
例如在加法群 \(\mathbb Z\) 上取点函数 \(\delta_0\)。乘法拉回在 \((m,n)\) 处为 \(\delta_{m+n,0}\)。若它来自代数张量积，就可写成某个有限和
\[
\delta_{m+n,0}=\sum_{i=1}^r a_i(m)b_i(n).
\]
固定 \(m\) 后，这个关于 \(n\) 的函数必落在 \(b_1,\ldots,b_r\) 的线性生成空间中。但左边随 \(m\) 变化给出所有点函数 \(\delta_{-m}\)，它们线性无关：在各自支集点求值即可逐个读出系数。这与有限维生成空间矛盾，故该两变量函数确实不在 \(k^{\mathbb Z}\otimes k^{\mathbb Z}\) 中。

\subsection{模与余模的区别}
\label{b5:subsec:2101:03}

有限群表示在基下写成 \(gv_j=\sum_i a_{ij}(g)v_i\)。与其对每个 \(g\) 分别保存这个矩阵，也可以把系数看成函数 \(a_{ij}\in k^G\)，一次写成 \(v_j\mapsto\sum_i v_i\otimes a_{ij}\)。群乘法在系数函数上的拉回，正好给出这种展开应满足的相容条件。这便是余模的来源。

\begin{definition}\label{b5:def:right-comodule}
设 \(k\) 为域，\((C,\Delta,\varepsilon)\) 为 \(k\) 余代数。\(k\) 向量空间 \(V\) 配 \(k\) 线性映射 \(\delta:V\to V\otimes_k C\)，使
\[
(\delta\otimes1)\delta=(1\otimes\Delta)\delta,\qquad
(1\otimes\varepsilon)\delta=1_V,
\]
时，称为右 \(C\)\term[yumo]{余模}[comodule]。第一式将两种继续展开余作用的次序联系起来，即下图交换：
\[
\begin{tikzcd}[column sep=large,row sep=large]
V\arrow[r,"\delta"]\arrow[d,"\delta"']
&V\otimes_kC\arrow[d,"1_V\otimes\Delta"]\\
V\otimes_kC\arrow[r,"\delta\otimes1_C"']
&V\otimes_kC\otimes_kC
\end{tikzcd}
\]
给定两个右 \(C\) 余模 \((V,\delta_V)\)、\((W,\delta_W)\)，若 \(k\) 线性映射 \(f:V\to W\) 使下图交换，就称它为余模同态：
\[
\begin{tikzcd}[column sep=large,row sep=large]
V\arrow[r,"\delta_V"]\arrow[d,"f"']
&V\otimes_kC\arrow[d,"f\otimes1_C"]\\
W\arrow[r,"\delta_W"']&W\otimes_kC
\end{tikzcd}
\]
\end{definition}
\symidx{\(\delta:V\to V\otimes_kC\)：右余模的余作用}
模是代数元素作用于向量的规则；余模是把一个向量展开成向量与系数函数的规则。只有适当的对偶条件成立时，两种语言才可直接转换。例如有限维代数 \(A\) 的左模等同于右 \(A^*\) 余模；取对偶基 \(a_i,a^i\)，余作用为
\(v\mapsto\sum_i a_iv\otimes a^i\)。
代数乘法公理与单位公理分别对应余结合与余单位。

\ProseParagraphBreak
一个二维例子已能看清余作用的含义。设 \(k\) 为域，\(C=k[t]\) 的余乘法为 \(\Delta(t)=t\otimes1+1\otimes t\)，余单位为 \(\varepsilon(t)=0\)。在 \(V=kv_1\oplus kv_2\) 上令
\[
\delta(v_1)=v_1\otimes1,\qquad
\delta(v_2)=v_2\otimes1+v_1\otimes t.
\]
在 \(v_2\) 上，余结合律两边都等于
\(v_2\otimes1\otimes1+v_1\otimes t\otimes1+v_1\otimes1\otimes t\)，余单位律也直接成立。把 \(t\) 代成 \(c\in k\)，便得到矩阵
\(\bigl(\begin{smallmatrix}1&c\\0&1\end{smallmatrix}\bigr)\)；余乘法中的两个加项对应参数相加。直线 \(kv_1\) 是子余模，商余模平凡，但将 \(v_2\) 换成 \(v_2+av_1\) 仍留下 \(v_1\otimes t\) 项，故这条子余模没有补。

\begin{proposition}\label{b5:prop:hopf-tensor-dual}
设 \(k\) 为域，\((H,\Delta,\varepsilon,S)\) 为 \(k\) 上的 Hopf 代数，\(V,W\) 为幺性左 \(H\) 模。对 \(h\in H\)，用 Sweedler 记号写 \(\Delta(h)=\sum h_{(1)}\otimes h_{(2)}\)。则张量积 \(V\otimes_kW\) 通过
\[
h(v\otimes w)=\sum h_{(1)}v\otimes h_{(2)}w
\]
成为左模，\(k\) 通过 \(\varepsilon\) 成为单位模。若 \(V\) 有限维，则
\[
(h\varphi)(v)=\varphi\bigl(S(h)v\bigr)
\]
给出 \(V^*\) 的左模结构。对 \(V\) 的基 \(e_1,\ldots,e_r\) 及对偶基 \(e_1^*,\ldots,e_r^*\)，求值与余求值映射
\[
\operatorname{ev}:V^*\otimes_kV\to k,\qquad
\varphi\otimes v\mapsto\varphi(v),
\]
\[
\operatorname{coev}:k\to V\otimes_kV^*,\qquad
1\mapsto\sum_i e_i\otimes e_i^*
\]
均为 \(H\) 模同态，并使 \(V^*\) 成为 \(V\) 的左对偶。
\end{proposition}
\begin{proof}
\(\Delta\) 保持乘法使两次作用等于乘积作用，余结合使三重张量的两种括号相容，余单位使 \(k\otimes V\to V\) 相容。对偶上依次作用 \(a,b\) 得到
\(\varphi\bigl(S(b)S(a)v\bigr)=\varphi\bigl(S(ab)v\bigr)\)，所以是左作用。求值映射的等变性为
\[
\sum\bigl(h_{(1)}\varphi\bigr)\bigl(h_{(2)}v\bigr)
=\varphi\left(\sum S\bigl(h_{(1)}\bigr)h_{(2)}v\right)
=\varepsilon(h)\varphi(v).
\]
为核验余求值，使用典范同构
\[
\Theta:V\otimes_kV^*\xrightarrow{\sim}\operatorname{End}_k(V),\qquad
\Theta(v\otimes\varphi)(x)=\varphi(x)v.
\]
张量 \(\operatorname{coev}(1)\) 对应 \(1_V\)，因而不依赖基的选择。对任意 \(x\in V\)，有
\[
\begin{aligned}
\Theta\bigl(h\operatorname{coev}(1)\bigr)(x)
&=\sum_{i,(h)} e_i^*\bigl(S(h_{(2)})x\bigr)h_{(1)}e_i\\
&=\sum_{(h)}h_{(1)}S(h_{(2)})x
=\varepsilon(h)x.
\end{aligned}
\]
由 \(\Theta\) 单射可知
\(h\operatorname{coev}(1)=\varepsilon(h)\operatorname{coev}(1)\)，故余求值也是 \(H\) 模同态。略去通常的结合与单位同构，两条蛇形复合分别将
\[
v\longmapsto\sum_i e_i\,e_i^*(v)=v,\qquad
\varphi\longmapsto\sum_i\varphi(e_i)e_i^*=\varphi.
\]
这就给出左对偶的两条恒等式；以上计算只用了对极公理。
\end{proof}
在群代数上，张量公式就是对角群作用；在包络代数上，它是 Lie 代数的 Leibniz 型作用。普通交换 \(v\otimes w\mapsto w\otimes v\) 是否等变，要看余乘法是否余交换；\subsecref{b5:subsec:B01:03}给出一个非余交换的例子。上述基础可对照\cite[Definitions 5.2.2, 5.3.2; Propositions 5.3.5--5.3.6; Corollary 5.3.7]{EtingofEtAl2015TensorCategories}。
